\section{Commuting forms with division coefficients}\label{sec:forms}

We state the construction theorem from the companion article
\cite{ArtinianEGH} in the precise form used below. A central
division algebra over \(k\) is a finite-dimensional
\(k\)-algebra with center \(k\) in which every nonzero element is invertible.
In \(\Delta[x]\), the variables \(x_i\) are central and have degree one,
and elements of \(\Delta\) have degree zero.

\begin{theorem}[Commuting-form construction, {\cite[Theorem~1.1]{ArtinianEGH}}]
\label{thm:forms}
Let \(f_1,\ldots,f_n\) be a homogeneous regular sequence in
\(\C[x_1,\ldots,x_n]\), of degrees \(2\le a_1\le\cdots\le a_n\).
There are a field extension \(k/\C\), a central division \(k\)-algebra
\(\Delta\), and a commutative graded \(k\)-subalgebra
\[
 k[x_1,\ldots,x_n]\subseteq A\subseteq N=\Delta[x_1,\ldots,x_n]
\]
generated by finitely many elements of degree one, containing
\(t_1,\ldots,t_n\in A_1\), such that
\begin{equation}\label{eq:triangular-forms}
 t_h^{a_h}=c_h f_h+\sum_{\ell<h}B_{h\ell}f_\ell
                  +\sum_{\ell>h}C_{h\ell}t_\ell
 \qquad(1\le h\le n).
\end{equation}
Here \(c_h\in k^\times\), \(B_{h\ell},C_{h\ell}\in A\), and every summand
has degree \(a_h\). Moreover,
\[
 \{t^\alpha=t_1^{\alpha_1}\cdots t_n^{\alpha_n}:
       \alpha\in\N^n\}
\]
is a graded left-\(\Delta\) basis of the \emph{full} algebra \(N\).
\end{theorem}

The basis assertion in Theorem~\ref{thm:forms} is stronger than a basis of the finite
complete-intersection quotient: every monomial \(t^\alpha\), without
an upper bound on its exponents, occurs in the decomposition of \(N\).
This is what will make \(N\) a free module over a polynomial ring and
allow us to transfer an entire resolution. The Hilbert-function
consequence of the foundation theorem alone would not suffice.

\begin{remark}\label{rem:coefficients}
Commutativity of the forms does not mean that their individual
\(\Delta\)-coefficients commute. The proof below always places those
coefficients on the left and multiplies by the \(t_i\) on the right.
The construction itself has no hypothesis on the additional generators
of \(I\).
\end{remark}
